\chapter[Q\&A: Transient Simulation]
{Questions and Answers for Transient Simulation}
\label{app:qa-transient-simulation}
\index{transient simulation!questions and answers}
\index{LTspice!troubleshooting}

This appendix collects practical questions that frequently arise when using
the transient-simulation toolbox. The answers refer to the conventions and
models used in this book.

\section{Basic Simulation Setup}

\begin{formulaBox}{Which parameters should be present on every schematic?}
The following global parameters are recommended:
\begin{lstlisting}[style=LTspice]
.param tdv=1 tdt=1
.param fn1=50 omega_n={2*pi*fn1}
.options vntol=1m abstol=1u
\end{lstlisting}

Use \texttt{fn1=60} for a \SI{60}{\hertz} system. The values
\texttt{vntol=1m} and \texttt{abstol=1u} correspond to
\SI{1}{\milli\volt} and \SI{1}{\micro\ampere}; they are included in the
toolbox template \texttt{template.asc}.

For a \SIrange{50}{60}{\hertz} system, a useful maximum transient time step is
between \SI{0.1}{\milli\second} and \SI{0.6}{\milli\second}. A smaller step
improves temporal resolution but increases simulation time. The alternate
solver is recommended under \textsc{Tools} $\rightarrow$
\textsc{Control Panel} $\rightarrow$ \textsc{SPICE}.

If a result is physically implausible, first check transformer orientation and
all engineering prefixes. In LTspice, \texttt{m} means milli; use
\texttt{Meg} for mega. A value of \texttt{400} means \SI{400}{\volt}, not
\SI{400}{\kilo\volt}.
\end{formulaBox}

\begin{formulaBox}{How can steady state be reached quickly?}
The voltage-source parameters \texttt{texp} and \texttt{tpwl} can be used to
ramp the source voltage. Values between approximately
\SI{10}{\milli\second} and \SI{100}{\milli\second} normally produce a final
voltage buildup over roughly \SIrange{50}{500}{\milli\second}.

Transformer energization can additionally be improved by coordinating the
residual-flux parameter \texttt{psi\_0} with the switching instant. For a
phase-R voltage zero crossing followed by a positive voltage half-cycle:
\begin{itemize}
\item \texttt{psi\_0=-1} can place the flux close to its steady-state path;
\item \texttt{psi\_0=0} can produce approximately twice rated peak flux;
\item \texttt{psi\_0=1} can produce approximately three times rated peak flux.
\end{itemize}
At a voltage maximum, \texttt{psi\_0=0} normally gives the smallest flux
offset. The exact polarity depends on the winding and voltage sign convention.
\end{formulaBox}

\begin{formulaBox}{What can be done if the circuit does not converge?}
Use the following sequence:
\begin{enumerate}
\item Verify units, polarities, initial conditions, floating nodes, and switch
      timing.
\item Reduce the maximum time step.
\item Use the alternate solver.
\item Provide realistic parasitic resistances and DC paths.
\item If necessary, relax \texttt{vntol}, \texttt{abstol}, or
      \texttt{reltol} moderately.
\end{enumerate}
Increasing tolerances can help convergence but reduces numerical accuracy;
repeat the final simulation with tighter tolerances as a check.
\end{formulaBox}

\begin{formulaBox}{How can subcircuit currents and voltages be accessed?}
Open
\begin{center}
\textsc{Tools} $\rightarrow$ \textsc{Control Panel} $\rightarrow$
\textsc{Save Defaults}
\end{center}
and enable \textbf{Save Subcircuit Node Voltages} and
\textbf{Save Subcircuit Device Currents}. Restart or rerun the simulation after
changing these settings.
\end{formulaBox}

\begin{formulaBox}{How are power, energy, and current direction displayed?}
Hold \textsc{Alt} and click a component to plot its instantaneous power, based
on the voltages and currents at its pins. Hold \textsc{Ctrl} and click the
power trace label to integrate the trace over the displayed interval and obtain
energy.

LTspice defines device-pin currents as positive into the pin. Resistors,
capacitors, and inductors display a current arrow. Mirroring a symbol also
mirrors this reference direction.
\end{formulaBox}

\begin{formulaBox}{What must be considered for LTspice inductors?}
LTspice inductors may contain a default series resistance of
\SI{1}{\milli\ohm}. Right-click the inductor and set \texttt{Rser=0} if an
ideal inductor is required. The global default can be changed under
\textsc{Tools} $\rightarrow$ \textsc{Control Panel} $\rightarrow$
\textsc{Hacks} by enabling the option that defaults inductors to
\texttt{Rser=0}.

Removing all loss can create an undamped resonant circuit and may worsen
convergence. Use a physically meaningful resistance whenever possible.
\end{formulaBox}

\section{Transformer Models and Magnetic Quantities}

\begin{formulaBox}{How is a nonlinear transformer changed to a linear model?}
Right-click the transformer and set
\begin{lstlisting}[style=LTspice]
m=1
\end{lstlisting}
in the value or parameter field. This removes the high-order nonlinear term in
the toolbox magnetization characteristic.
\end{formulaBox}

\begin{formulaBox}{Which parameters determine the magnetization curve?}
The characteristic is controlled primarily by \texttt{a} and \texttt{m}:
\begin{equation}
i_\mu(\Phi)=I_{\mu N}
\left[
a\frac{\Phi}{\Phi_N}
+(1-a)\left(\frac{\Phi}{\Phi_N}\right)^m
\right].
\label{eq:qa-magnetization-curve}
\end{equation}

The parameter $a$ determines the linear share and the small-signal slope,
while $m$ controls the sharpness of saturation. Typical large-transformer
examples use $m=7$ or $m=9$. 

The quantities $I_{\mu N}$ and $\Phi_N$ must use the same amplitude convention.
If peak flux is used, the rated peak values are
\begin{equation}
\Phi_N=\frac{\sqrt{2}\,U_{N,\mathrm{ph}}}{\omega N},
\qquad
I_{\mu N}=\sqrt{2}\,I_{\mu N,\mathrm{rms}}.
\end{equation}
\end{formulaBox}

\begin{formulaBox}{How is residual flux specified?}
Use the transformer parameter \texttt{psi\_0} for a predefined residual-flux
pattern. Where supported, phase-specific parameters should be used for an
arbitrary residual-flux vector. The initialization time must match the desired
switching event.
\end{formulaBox}

\begin{formulaBox}{How is residual flux applied at an arbitrary time?}
The model briefly connects an initialization voltage source to the
\SI{1}{\farad} flux-integrator capacitor. The switch on-resistance and duration
determine the charging accuracy:
\begin{equation}
\tau=R_{\mathrm{on}}C.
\end{equation}
For example,
\begin{equation}
R_{\mathrm{on}}=10^{-5}\,\Omega,\quad C=\SI{1}{\farad}
\quad\Longrightarrow\quad
\tau=\SI{0.01}{\milli\second}.
\end{equation}
A closing duration of \SI{0.1}{\milli\second} corresponds to $10\tau$ and
charges the capacitor essentially to its final value. The rapid initialization
can produce a narrow voltage or current spike at the charging instant.
\end{formulaBox}

\begin{formulaBox}{Where can introductory LTspice tutorials be found?}
An introductory tutorial collection is available at
\url{http://www.simonbramble.co.uk/lt_spice/ltspice_lt_spice.htm}.
\end{formulaBox}

\begin{formulaBox}{How is the JUSTUS transformer model composed?}
The model combines an electrical circuit with a magnetic equivalent circuit.
Iron loss, leakage, and core paths can therefore be represented with magnetic
quantities. Under the magnetic-electrical duality used in the model:
\begin{itemize}
\item electrical series connections correspond to magnetic parallel paths;
\item electrical parallel connections correspond to magnetic series paths;
\item resistance corresponds to conductance under the chosen analogy;
\item voltage-source and current-source roles are interchanged;
\item T and $\Pi$ topologies are dual representations.
\end{itemize}

E.g to convert a inductance to a reluctance or a resistance to a reluctance for a winding with $N$ turns,
\begin{align}
\Rmag_\sigma&=\frac{N^2}{L_\sigma},\\
\Rmag_{R\mathrm{Fe}}&=\frac{N^2}{R_{\mathrm{Fe}}},
\end{align}
according to the scaling convention of the toolbox model. The placement of the
corresponding branches follows the electrical-magnetic duality described above.
\end{formulaBox}

\begin{formulaBox}{How can approximate data for a \SI{300}{\mega\volt\ampere} transformer be estimated?}
Assume a \SI{400}{\kilo\volt}/\SI{120}{\kilo\volt} Yy transformer,
$N_1=1000$, magnetic path length $l_m=\SI{10}{\metre}$, rated flux density
$B_N=\SI{1.8}{\tesla}$, and RMS magnetizing current
$I_{\mu N}=\SI{0.5}{\ampere}$.

Using RMS sinusoidal quantities,
\begin{align}
\Psi_{1N,\mathrm{rms}}
 &=\frac{U_{1N,LL}}{\sqrt{3}\,\omega}
  =\SI{735.1}{\weber},\\
\Phi_{N,\mathrm{rms}}
 &=\frac{\Psi_{1N,\mathrm{rms}}}{N_1}
  =\SI{0.7351}{\weber},\\
A_{\mathrm{Fe}}
 &=\frac{\Phi_{N,\mathrm{rms}}}{B_N}
  \approx\SI{0.408}{\metre\squared},\\
\Theta_{\mu N}&=N_1I_{\mu N}=\SI{500}{\ampere},\\
H_N&=\frac{\Theta_{\mu N}}{l_m}=\SI{50}{\ampere\per\metre},\\
\mu_r&=\frac{B_N}{\mu_0H_N}\approx2.86\times10^4,\\
\Rmag_m&=\frac{\Theta_{\mu N}}{\Phi_{N,\mathrm{rms}}}
 \approx\SI{680}{\per\henry},\\
L_{1h}&=\frac{N_1^2}{\Rmag_m}\approx\SI{1470}{\henry}.
\end{align}

The equivalent peak convention gives
$\widehat\Phi_N=\sqrt{2}\Phi_{N,\mathrm{rms}}
=\SI{1.0396}{\weber}$ and
$\widehat I_{\mu N}=\sqrt{2}I_{\mu N}$, leaving reluctance and inductance
unchanged. RMS and peak quantities must not be mixed in one calculation.
\end{formulaBox}

\begin{formulaBox}{What does the nonlinear curve predict at one, two, and three times rated flux?}
For $a=0.75$, $m=9$, and a reference current of
$I_{\mu N}=\SI{0.5}{\ampere}$ in a consistent amplitude convention,
\begin{align}
i_\mu(1\Phi_N)&=\SI{0.5}{\ampere},\\
i_\mu(2\Phi_N)&=\SI{64.75}{\ampere},\\
i_\mu(3\Phi_N)&=\SI{2461.5}{\ampere}.
\end{align}

For $l_m=\SI{10}{\metre}$ and $B_N=\SI{1.8}{\tesla}$, the corresponding
approximate magnetic quantities are
\begin{table}[H]
\centering
\renewcommand{\arraystretch}{1.12}
\begin{tabular}{ccccc}
\toprule
$\Phi/\Phi_N$ & $B$ & $H$ & $\mu_r$ & $L_{1h}$ \\
\midrule
1 & \SI{1.8}{\tesla} & \SI{50}{\ampere\per\metre} & $2.86\times10^4$ & \SI{1470}{\henry} \\
2 & \SI{3.6}{\tesla} & \SI{6475}{\ampere\per\metre} & 442 & \SI{22.7}{\henry} \\
3 & \SI{5.4}{\tesla} & \SI{2.46e5}{\ampere\per\metre} & 17.5 & \SI{0.90}{\henry} \\
\bottomrule
\end{tabular}
\end{table}

These are point estimates based on the assumed algebraic curve. In a dynamic
simulation, winding voltage drop, leakage flux, and waveform distortion reduce
the effective flux and RMS current. Even at three times rated flux, the
estimated main inductance remains above the example leakage inductance
$L_\sigma\approx\SI{0.1697}{\henry}$.
\end{formulaBox}

\begin{formulaBox}{How can tank reluctance be estimated?}
The tank path is strongly geometry- and frequency-dependent. A placeholder
inductance of a few henries, or roughly 10 to 100 times the winding leakage
inductance, may be useful for sensitivity studies, but it is not a validated
transformer parameter. Measurements or field calculations are required for an
accurate value.
\end{formulaBox}

\begin{formulaBox}{How can transformer voltage and phase be regulated?}
The rudimentary tap parameter \texttt{l\_} changes the voltage in steps. For
example, \texttt{l\_=3} and \texttt{stproz=1.5} correspond to a nominal
\SI{4.5}{\percent} change, subject to the model sign convention.

The models \texttt{trafo\_tap\_l} and \texttt{trafo\_tap\_s} provide voltage
and phase-angle regulation. Versions with fixed parameters and versions with
control pins are available; the latter can be driven by a PWL source during a
transient simulation.
\end{formulaBox}

\section{Selecting and Analysing Models}

\begin{formulaBox}{How are alternative models selected?}
Right-click the element and open the \textbf{SpiceModel} drop-down list. It
contains models from the associated \texttt{.lib} file. The model
pin count and pin order must match the symbol! 
\end{formulaBox}

\begin{formulaBox}{How is a transient model converted to AC analysis?}
Replace transient sources by sources with AC magnitude and phase definitions,
and replace the \texttt{.tran} directive by an \texttt{.ac} directive. AC
analysis linearizes the circuit about its operating point; it is generally not
suitable for studying switching events or strongly nonlinear large-signal
behaviour.
\end{formulaBox}

\begin{formulaBox}{How can reactive-power transfer be recognized?}
Reactive-power transfer produces a phase displacement between the fundamental
voltage and current. For sinusoidal RMS phasors,
\begin{equation}
Q=\sqrt{3}\,U_{LL}I_L\sin\varphi.
\end{equation}
For distorted waveforms, calculate average active power from $u(t)i(t)$ and
use a harmonic-aware definition of reactive power.
\end{formulaBox}

\begin{formulaBox}{Why do many elements have control pins?}
Control pins allow parameters such as voltage level, tap position, phase shift,
switch state, or model coefficient to change during the simulation. A PWL
source is commonly used to drive a discrete or continuous control input.
\end{formulaBox}

\section{Inductive Coupling and Flux}

\begin{formulaBox}{How is leakage between two coupled inductors calculated?}
For self-inductances $L_1,L_2$ and mutual inductance $M$,
\begin{align}
k&=\frac{M}{\sqrt{L_1L_2}},\\
\sigma&=1-\frac{M^2}{L_1L_2}=1-k^2.
\end{align}
If the individual leakage factors are $\sigma_1$ and $\sigma_2$,
\begin{equation}
\sigma=\sigma_1+\sigma_2-\sigma_1\sigma_2.
\end{equation}
For a symmetric transformer, $\sigma_1=\sigma_2$, hence
\begin{equation}
\sigma=2\sigma_1-\sigma_1^2,
\qquad k=1-\sigma_1,
\qquad L_{1\sigma}=(1-k)L_1.
\end{equation}

With $L_{1\sigma}=\SI{0.1697}{\henry}$ and
$L_1=\SI{1470}{\henry}$,
\begin{align}
\sigma_1&=1.1544\times10^{-4},\\
\sigma&=2.3087\times10^{-4},\\
k&=\boxed{0.9998846}.
\end{align}
The voltage ratio and inductances satisfy
\begin{equation}
\frac{N_1}{N_2}=\sqrt{\frac{L_1}{L_2}}
\end{equation}
when both inductances use the same magnetic path and reference convention.
\end{formulaBox}

\begin{formulaBox}{What is the rated peak flux for 1000 turns?}
For sinusoidal phase voltage $U_{N,\mathrm{ph}}$,
\begin{equation}
\widehat\Psi_N=\frac{\sqrt{2}\,U_{N,\mathrm{ph}}}{\omega},
\qquad
\widehat\Phi_N=\frac{\widehat\Psi_N}{N}.
\end{equation}

At \SI{50}{\hertz} and $N=1000$:
\begin{table}[H]
\centering
\renewcommand{\arraystretch}{1.12}
\begin{tabular}{lcc}
\toprule
Voltage applied to winding & Connection interpretation & $\widehat\Phi_N$ \\
\midrule
\(\SI{10}{\kilo\volt}/\sqrt{3}\) & \SI{10}{\kilo\volt} star phase & \SI{0.0260}{\weber} \\
\SI{10}{\kilo\volt} & Winding voltage & \SI{0.0450}{\weber} \\
\(\SI{20}{\kilo\volt}/\sqrt{3}\) & \SI{20}{\kilo\volt} star phase & \SI{0.0520}{\weber} \\
\SI{20}{\kilo\volt} & Winding voltage & \SI{0.0900}{\weber} \\
\(\SI{110}{\kilo\volt}/\sqrt{3}\) & \SI{110}{\kilo\volt} star phase & \SI{0.286}{\weber} \\
\SI{110}{\kilo\volt} & Winding voltage & \SI{0.495}{\weber} \\
\(\SI{220}{\kilo\volt}/\sqrt{3}\) & \SI{220}{\kilo\volt} star phase & \SI{0.572}{\weber} \\
\SI{220}{\kilo\volt} & Winding voltage & \SI{0.990}{\weber} \\
\(\SI{400}{\kilo\volt}/\sqrt{3}\) & \SI{400}{\kilo\volt} star phase & \SI{1.0396}{\weber} \\
\bottomrule
\end{tabular}
\end{table}
\end{formulaBox}

\begin{formulaBox}{What are the time constants of RL and RC circuits?}
\begin{equation}
\tau_{RL}=\frac{L}{R},
\qquad
\tau_{RC}=RC.
\end{equation}
A first-order response reaches approximately \SI{99.3}{\percent} of its final
change after $5\tau$ and \SI{99.995}{\percent} after $10\tau$.
\end{formulaBox}

\begin{formulaBox}{What do \texttt{tpsid} and \texttt{ron} mean?}
The parameter \texttt{tpsid} specifies how long the initialization source is
connected to the \SI{1}{\farad} flux-integrator capacitor. The parameter
\texttt{ron} is the switch on-resistance. Suitable example values are
\begin{lstlisting}[style=LTspice]
tpsid=0.1m  ron=1e-5
\end{lstlisting}
giving $\tau=R_{\mathrm{on}}C=\SI{0.01}{\milli\second}$ and an initialization
interval of $10\tau$.
\end{formulaBox}

\begin{formulaBox}{How can the magnetization curve be plotted?}
Open the transformer secondary. In the waveform viewer, change the horizontal
axis from time to primary current and plot internal flux or flux linkage versus
that current. A nonzero iron-loss parameter \texttt{p0} produces a loop-like
separation because the loss current is superimposed; it should not be confused
with a fully physical hysteresis model unless the model explicitly contains
hysteresis state equations.
\end{formulaBox}

\section{Lines, Cables, and Compensation}

\begin{formulaBox}{How is cable charging power calculated?}
For line-to-line RMS voltage $U_N$, phase-to-earth capacitance $C_E$, and
phase-to-phase capacitance $C_L$ per phase set,
\begin{align}
Q_0&=Q_{CE}+Q_{CL}
   =U_N^2\omega C_E+3U_N^2\omega C_L,\\
I_0&=\frac{Q_0}{\sqrt{3}U_N}.
\end{align}
For a shielded single-core cable, $C_L$ is commonly negligible and
$Q_0\approx Q_{CE}$. The capacitive earth-fault current is then approximately
\begin{equation}
I_E=3I_0=\frac{3Q_{CE}}{\sqrt{3}U_N}.
\end{equation}
\end{formulaBox}

\begin{formulaBox}{How are Petersen coils and shunt reactors sized?}
For a Petersen coil connected from the transformer neutral to earth,
\begin{align}
X_L&=\frac{U_N/\sqrt{3}}{I_N}, &
L&=\frac{X_L}{\omega}.
\end{align}
For an example with $U_N=\SI{110}{\kilo\volt}$ and
$I_N=\SI{250}{\ampere}$, the coil loss may initially be estimated with
$R\approx X_L/100$ and later replaced by measured data.

For a three-phase reactor with total rating $Q_N$:
\begin{align}
L_\Delta&=\frac{3U_N^2}{\omega Q_N},\\
L_Y&=\frac{U_N^2}{\omega Q_N}.
\end{align}
Each delta branch sees $U_N$, whereas each star branch sees $U_N/\sqrt{3}$.

If $C_E=\SI{1000}{\nano\farad}$ at \SI{110}{\kilo\volt},
\begin{align}
Q_0&=U_N^2\omega C_E\approx\SI{3.80}{\mega\varunit},\\
I_0&=\frac{Q_0}{\sqrt{3}U_N}\approx\SI{19.95}{\ampere},\\
I_E&=3I_0\approx\SI{59.85}{\ampere}.
\end{align}
A Petersen coil rated approximately \SI{110}{\kilo\volt} and
\SI{60}{\ampere} therefore provides the required compensation.
\end{formulaBox}

\begin{formulaBox}{How does coupling influence parallel lines?}
Lines on the same route have mutual inductive and capacitive coupling. Depending
on current direction and phase arrangement, mutual induction can reduce or
increase the effective series inductance. In the idealized same-direction case,
$k\rightarrow1$ makes the differential leakage inductance approach zero.

The toolbox contains coupled-line models for one, two, and four circuits. Their
coupling matrices were calculated with ATP/ATPDraw and converted to LTspice
format with the supplied Excel conversion file.
\end{formulaBox}

\begin{formulaBox}{How is source impedance obtained from short-circuit power?}
For a balanced three-phase network,
\begin{equation}
Z_N=\frac{U_N^2}{S_k}.
\end{equation}
Examples are
\begin{align}
U_N=\SI{400}{\kilo\volt},\quad S_k=\SI{40}{\giga\volt\ampere}
&\Rightarrow Z_N=\SI{4.0}{\ohm},\\
U_N=\SI{110}{\kilo\volt},\quad S_k=\SI{7.5}{\giga\volt\ampere}
&\Rightarrow Z_N=\SI{1.613}{\ohm}.
\end{align}
If no resistance information is available, the first approximation is an
inductive source impedance $jX_N$ with $X_N\approx Z_N$.
\end{formulaBox}

\begin{formulaBox}{How is the total capacitance of an overhead line calculated?}
With phase-to-earth capacitance $C_E$ and phase-to-phase capacitance $C_L$, the
effective positive-sequence charging capacitance is
\begin{equation}
C_{\mathrm{sum}}=C_E+3C_L.
\end{equation}
The capacitive reactive power is
\begin{equation}
Q_C=U_N^2\omega C_{\mathrm{sum}}
   =\frac{U_N^2}{X_C}.
\end{equation}
\end{formulaBox}

\begin{formulaBox}{How is the reactive power of a shunt reactor calculated?}
For a balanced three-phase reactor connected in star,
\begin{equation}
Q_L=\frac{U_N^2}{X_L}
   =\frac{U_N^2}{\omega L},
\end{equation}
where $L$ is the inductance of each phase branch.
\end{formulaBox}

\begin{formulaBox}{What is the difference between star and delta loads with identical branch impedances?}
If the same branch resistance or impedance is used in both connections at the
same line-to-line voltage, each delta branch receives $\sqrt{3}$ times the star
branch voltage. Therefore,
\begin{equation}
P_\Delta=3P_Y,
\qquad Q_\Delta=3Q_Y.
\end{equation}
\end{formulaBox}

\section{Plot Files and Verification}

\begin{formulaBox}{How is a saved plot configuration displayed?}
Many examples include saved plot-setting files. In the waveform viewer, use
\textsc{Plot Settings} $\rightarrow$ \textsc{Open Plot Settings} and select
the required file. If several configurations are supplied, open them one at a
time to reproduce the intended trace groups.
\end{formulaBox}

\begin{formulaBox}{How should a transient simulation be verified?}
Always repeat the simulation with a smaller maximum time step. A credible
result should change only slightly as the time step is reduced. Also check:
\begin{itemize}
\item energy and power balance;
\item expected steady-state RMS values and phase sequence;
\item sensitivity to tolerances and parasitic resistances;
\item agreement between transient, AC, and hand calculations where applicable;
\item correct current reference directions and voltage bases;
\item plausibility of resonant frequencies and decay time constants.
\end{itemize}
\end{formulaBox}

\begin{chaptersummary}[Appendix Summary]
Reliable transient simulation begins with consistent global parameters,
correct units, suitable time steps, and physically meaningful damping. Source
ramps and coordinated residual flux reduce initialization transients, while
the alternate solver and carefully selected tolerances improve numerical
robustness. Every important result should be repeated with a
smaller time step and compared with analytical estimates.
\end{chaptersummary}
